\documentclass[12pt]{article}

\addtolength{\textwidth}{.8in}
\addtolength{\hoffset}{-.4in}
\addtolength{\textheight}{1in}
\addtolength{\voffset}{-1in}

\usepackage{wasysym}
\usepackage{multicol}
\usepackage{amsfonts}
\usepackage{graphicx}
\usepackage{hyperref}
\usepackage{enumerate}
\usepackage{parskip}
\usepackage{paralist}
\usepackage{upquote} % to convert funny quotes to straight quotes
\setlength{\parindent}{15pt}
\pagestyle{empty}

\renewcommand*{\thefootnote}{\fnsymbol{footnote}}

\begin{document}

\begin{center}
\begin{Large}
DATA 219 Notes -- Jan. 25, 2018
\end{Large}
\end{center}

\section{Administuff}

\begin{compactitem}

\item Ink it in: \textbf{DataFest: Penn State, April 13-15}! Details
forthcoming. I'll need a firm yes-or-no from everyone by class time on
\textbf{Thursday February 8th.}

\end{compactitem}

\section{Generating synthetic data sets}

One very useful skill is the ability to quickly create \textit{synthetic}
(artificial; generated by your own code and random-number generators) data
sets that have certain properties. Sometimes we use such data to sanity check
the results of our code with ``idealistic" (simplified and known) inputs.
Sometimes we use synthetic data as a baseline with which to compare real-world
data sets that we suspect have similar characteristics. And sometimes we
simply don't have access to relevant real-world data but we need inputs into
some simulation process.

\subsection{Uncorrelated numeric variables}

We learned last time how to use NumPy's random number generation to create
arrays according to several statistical distributions. Suppose you wanted to
create a data set with two \textit{uncorrelated} variables. You'd just
generate two arrays, one at a time:

\begin{Verbatim}[fontsize=small,samepage=true,frame=none]
heights = np.random.normal(65, 8, 100)
sat_scores = np.random.normal(1400, 350, 100)
\end{Verbatim}

Scatterplotting these, or running a Pearson correlation, should reveal they
are independent of each other.

\subsection{Totally correlated numeric variables}

If two variables are \textit{completely} correlated, then one is simply a
direct (linear) function of the other. That's easy to accomplish by creating
only \textit{one} randomly, and then using it to compute the other:

\begin{Verbatim}[fontsize=small,samepage=true,frame=none]
weights_lbs = np.random.normal(150, 30, 100)
weights_kg = weights_lbs / 2.2
\end{Verbatim}

Scatterplotting these should give a perfectly straight line, and Pearson
should tell us they're totally correlated (r=1).

\subsection{Somewhat correlated numeric variables}

The most common way to generate a data set in which two numerical variables
are somewhat, but not completely, correlated, is to generate one randomly, and
then to compute the second as a function of the first \textit{with added
noise}. ``Noise" is a somewhat misleading word that's used in a variety of
statistical situations: here it means ``other (linearly additive) influences
that affect the second variable besides just the value of the first." To wit:

\begin{Verbatim}[fontsize=small,samepage=true,frame=none]
sat_scores = np.random.normal(1400, 350, 100)
gpas = sat_scores * (2.5 / 1400) + np.random.normal(0, 1, 100)
\end{Verbatim}

Notice that the second component of \texttt{gpas} is a normally distributed
``noise" component \textit{with zero mean}. 

\subsection{Generating \textit{categorical} variables}

NumPy provides a way to generate a categorical data set with certain
frequencies, through its \texttt{np.random.choice()} function:

\begin{Verbatim}[fontsize=small,samepage=true,frame=none]
np.random.choice(['Steelers','Redskins','Giants'])
'Redskins'
\end{Verbatim}

Useful arguments include:

\begin{compactitem}
\item \texttt{size} -- the number of elements you want in the resulting array.
\item \texttt{replace} -- if \texttt{False}, the elements will be drawn from
your list \textit{without replacement} (meaning once a value is chosen, it
will not be chosen again). This is often useful.
\item \texttt{p} -- a list of \textit{probabilities} (in the range 0 to 1)
specifying how likely each element in the list is to be chosen. (Note these
must sum to 1, and NumPy dumbly will not normalize it for you.)
\end{compactitem}

A more complicated example:
\begin{Verbatim}[fontsize=small,samepage=true,frame=none]
np.random.choice(['Steelers','Redskins','Giants'], size=10, p=[.3,.6,.1])
['Redskins' 'Redskins' 'Giants' 'Redskins' 'Redskins' 'Steelers' 'Redskins'
 'Steelers' 'Redskins' 'Giants']
\end{Verbatim}

\section{JSON}

JSON (pronounced like the name ``Jason") is a simple, human-readable text file
format that is very commonly used for information storage and interchange. (It
stands for "JavaScript Object Notation," but don't be fooled: it doesn't have
much to do with the JavaScript language except historically. \textit{Any}
language, including Python, can read/write JSON data.)

The really nice things are (1) it's simple, and easy on the eyes (unlike
binary formats, or XML), and (2) it structures information hierarchically in
list-ish, dictionary-ish, string-ish type elements that will be instantly
familiar to you. A JSON ``document" is, in fact, \textit{essentially a Python
list or dictionary with embedded contents.}

Here's an example file in JSON format:

\small
\begin{Verbatim}[fontsize=small,samepage=true,frame=none]
{"name": "Luke Skywalker", "rank": "Jedi", "lightsaber color": "blue",
"episodes": ["IV", "V", "VI", "VII", "VIII"], "relations": {"Leia": "sister",
"Han": "friend", "Artoo": "servant", "Rey": "student"}}
\end{Verbatim}
\normalsize

Its structure is instantly apparent: a dictionary whose values are strings,
lists, and other dictionaries, as appropriate.

Probably the best thing about JSON is that once we read and parse a JSON file
or other information source, we have a bona fide Python object that we can
treat in the ordinary way.

In Python, we have the \texttt{json} package, which you can import like any
other:

\begin{Verbatim}[fontsize=small,samepage=true,frame=none]
    import json
\end{Verbatim}

There are four functions that matter:
\small
\begin{compactitem}
\item \texttt{json.dump}(\textit{someObject}, \textit{aFile}) - Dump the Python
object to a file in JSON text format.
\item \texttt{json.dumps}(\textit{someObject}) - Dump the Python
object to a string in JSON text format (and return it).
\item \texttt{json.load}(\textit{aFile}) - Build and return a Python
object from the contents of the JSON text file.
\item \texttt{json.loads}(\textit{someString}) - Build and return a Python object
from the JSON text string.
\end{compactitem}

\normalsize
So, for instance, to load a file \texttt{info.json} into a Python object
called ``\textit{my\_stuff}":

\small
\begin{Verbatim}[fontsize=small,samepage=true,frame=none]
    with open('info.json','r') as f:
        my_stuff = json.load(f)
\end{Verbatim}

\normalsize
We could parse and access the Luke Skywalker file just like this:

\small
\begin{Verbatim}[fontsize=small,samepage=true,frame=none]
    with open('luke.json','r') as f:
        luke = json.load(f)
    print(luke['lightsaber color'])
    blue

    print(luke['episodes'][2])
    VI
\end{Verbatim}

\normalsize
Simple as that!

\end{document}
